\section*{摘要}

\noindent
针对车载平衡滚球运动控制任务，设计了一套由车辆循迹、钢球视觉测量和摆杆闭环控制组成的分布式系统。车辆侧以CH32V307为主控，采用八路光电加权误差、循迹PD外环和双轮PI速度内环，实现差速循迹、起终点识别与主动停车；K230在灰度感兴趣区域内完成圆目标检测、运动预测选通和位置滤波，并通过串口向STM32发送球位；STM32采用位置比例、积分和球速反馈相结合的控制器驱动舵机，同时接收车辆侧\SI{50}{Hz}运动状态帧，叠加纵向加速度与偏航角速度前馈。最终标定的机械中心为314像素，$-5\,\si{cm}$与$+5\,\si{cm}$目标分别对应190像素与445像素。

保留的一次整车工程试验以506像素（约$+7.4\,\si{cm}$）为目标，整圈用时\SI{28.00}{s}；按\SI{25}{px/cm}的保守比例换算，1400个有效控制样本的平均绝对误差为\SI{0.78}{cm}，95百分位误差为\SI{1.52}{cm}，最大绝对误差为\SI{1.76}{cm}，全程未出现丢球保护切换。试验表明系统已完成偏置目标下的整车闭环，但弯道阶段仍是误差主要来源，尚不能据此判定全程\SI{1}{cm}精度达标。无线图传的软件链路已完成双通道采集与MJPEG发送设计，因未保留场外接收、完整录像及回放记录，本文不将其计为已通过项目。

\noindent\textbf{关键词：}循迹小车；机器视觉；滚球控制；级联控制；运动前馈
